\subsection{Vertical Translation} 
To translate a point upward by $c$, we add $c$ to its $y$-coordinate. To translate it downward by $c$, we subtract $c$ from its $y$-coordinate.
\begin{Example}
In the figure below, the graph of:
\begin{itemize}
\item $h(x)$ is obtained by translating the graph of $f(x)$ three units upward
\item $k(x)$ is obtained by translating the graph of $f(x)$ two units downward
\end{itemize}
Give expressions for $h(x)$ and $k(x)$ in terms of $f(x)$.
\begin{icpic}[x=0.5\linespace,y=0.5\linespace]
\useasboundingbox (-7,-7) rectangle (7,7);
\setgraph{7}{7}{7}{7}
\GraphSmall
\foreach \y/\lbl/\col in {0/f/colf,3/h/colg,-2/k/colh} {
	\draw[thick,\col] (-4,1+\y) -- (-1,3+\y) -- (3,-3+\y) -- (6,-1+\y);
	\draw[\col] (5,-1+\y) node[anchor=north west] {$y=\lbl(x)$};
}
\end{icpic}
\funtable{4}{f}{h}{k}
\end{Example}
%%%%%%%%%%%%%%%%%%%%%%%%%%%%%%%%%%%%%%%%%%%%%%%%%%%%%%%%%%
\begin{displaybox}[Vertical Translation]
For any number $c>0$:
\begin{itemize}
\item If $g(x)=f(x)+c$, the graph of $g$ is the same as the graph of $f$, translated to the \makeblank[1in] by $c$ units.
\item If $g(x)=f(x)-c$, the graph of $g$ is the same as the graph of $f$, translated to the \makeblank[1in] by $c$ units.
\end{itemize}
\end{displaybox}
\begin{Note}
The direction of translation is the \makeblank[1.5in] as the sign of the number.
\end{Note}